% !TeX program = xelatex
% !TeX encoding = UTF-8
% Test Plan -- single self-contained file (preamble, metadata, cover, body).
% Build from this directory:  latexmk -xelatex main.tex   (or xelatex main.tex twice)
\documentclass[10pt,letterpaper]{article}

% -------------------- Metadata --------------------
\newcommand{\ProjectName}{MEDICATION SAFETY ASSISTANT}
\newcommand{\GroupName}{Medication Safety Assistant Team}
\newcommand{\GroupNumber}{Group 02}
\newcommand{\OrganizationName}{University of Science, VNU-HCM}
\newcommand{\DocumentAuthors}{Cao Hoàng Long, Hoàng Yến Trang, Lỗ Chí Thành}
\newcommand{\DocumentTitle}{Test Plan}
\newcommand{\DocumentIdentifier}{PA4-Group02-TP}
\newcommand{\DocumentVersion}{1.0}
\newcommand{\DocumentDate}{28 November 2026}
\newcommand{\ConfidentialityText}{Confidential}

% -------------------- Fonts and layout (same as PA3 documents) --------------------
\usepackage{fontspec}
\defaultfontfeatures{Ligatures=TeX}
\IfFontExistsTF{Times New Roman}{\setmainfont{Times New Roman}}{\setmainfont{Liberation Serif}}
\IfFontExistsTF{Arial}{\setsansfont{Arial}}{\setsansfont{Liberation Sans}}
\usepackage[letterpaper,margin=1in,headheight=34pt,headsep=16pt,footskip=28pt]{geometry}
\usepackage{setspace}\setstretch{1.08}
\usepackage{microtype}
\usepackage{array,tabularx,longtable,booktabs}
\usepackage{graphicx}
\usepackage{float}
\usepackage{caption}
\usepackage{enumitem}
\usepackage{titlesec}
\usepackage{fancyhdr}
\usepackage{lastpage}
\usepackage{xcolor}
\usepackage{tocloft}
\usepackage[unicode,hidelinks]{hyperref}

\setlength{\parindent}{0pt}
\setlength{\parskip}{5pt}
\setlength{\tabcolsep}{5pt}
\renewcommand{\arraystretch}{1.23}
\emergencystretch=2em
\setcounter{tocdepth}{2}
\setlist[itemize]{topsep=3pt,itemsep=2pt,parsep=0pt}
\renewcommand{\contentsname}{Table of Contents}
\titleformat{\section}{\normalfont\bfseries\fontsize{12}{15}\selectfont}{\thesection.}{0.6em}{}
\titleformat{\subsection}{\normalfont\bfseries\normalsize}{\thesubsection}{0.6em}{}
\titlespacing*{\section}{0pt}{8pt}{3pt}
\titlespacing*{\subsection}{0pt}{6pt}{3pt}
\renewcommand{\cfttoctitlefont}{\hfill\normalfont\bfseries\fontsize{17}{21}\selectfont}
\renewcommand{\cftaftertoctitle}{\hfill\mbox{}}
\renewcommand{\cftsecleader}{\cftdotfill{\cftdotsep}}
\renewcommand{\cftsecaftersnum}{.}
\tocloftpagestyle{fancy}
\newcolumntype{L}[1]{>{\raggedright\arraybackslash}p{#1}}
\newcolumntype{Y}{>{\raggedright\arraybackslash}X}

% -------- PA4 update marking: text added or changed in PA4 is blue --------
\definecolor{pafourblue}{RGB}{0,82,204}
\newcommand{\upd}[1]{{\color{pafourblue}#1}}
\newenvironment{updated}{\par\color{pafourblue}}{\par}
\newcommand{\frontmattertitle}[1]{\begin{center}{\normalfont\bfseries\fontsize{16}{20}\selectfont #1\par}\end{center}}

\pagestyle{fancy}
\fancyhf{}
\fancyhead[L]{\begin{minipage}[b]{.73\textwidth}\raggedright\fontsize{9}{11}\selectfont
  \ProjectName\\[-1pt]{\fontsize{8}{10}\selectfont\DocumentTitle\enspace(\DocumentIdentifier)}\end{minipage}}
\fancyhead[R]{\begin{minipage}[b]{.24\textwidth}\raggedleft\fontsize{9}{11}\selectfont
  Version \DocumentVersion\\[-1pt]{\fontsize{8}{10}\selectfont\DocumentDate}\end{minipage}}
\fancyfoot[L]{\fontsize{8.5}{10}\selectfont\OrganizationName{} -- Course Project}
\fancyfoot[C]{\fontsize{8}{10}\selectfont\ConfidentialityText}
\fancyfoot[R]{\fontsize{8.5}{10}\selectfont Page \thepage\ of \pageref{LastPage}}
\renewcommand{\headrulewidth}{0.45pt}
\renewcommand{\footrulewidth}{0.45pt}
\fancypagestyle{cover}{\fancyhf{}\renewcommand{\headrulewidth}{0pt}\renewcommand{\footrulewidth}{0pt}}
\hypersetup{pdftitle={\DocumentTitle\ - \ProjectName},pdfauthor={\DocumentAuthors}}

\begin{document}

% ------------------ Cover ------------------
\thispagestyle{cover}
\begin{center}
  \vspace*{.35in}
  {\bfseries\fontsize{13}{17}\selectfont \ProjectName\par}
  \vspace*{1.50in}
  {\bfseries\fontsize{22}{27}\selectfont \DocumentTitle\par}
  \vspace{8pt}
  {\fontsize{12}{16}\selectfont (System Testing)\par}
  \vspace{.92in}
  {\fontsize{12}{15}\selectfont Version \DocumentVersion\par}
  \vfill
  {\fontsize{10}{14}\selectfont
    \begin{tabular}{@{}r@{\quad}l@{}}
      \textbf{Document identifier:} & \DocumentIdentifier\\
      \textbf{Prepared by:} & \GroupNumber{} -- \GroupName\\
      \textbf{Authors:} & \DocumentAuthors\\
      \textbf{Organization:} & \OrganizationName\\
      \textbf{Date:} & \DocumentDate
    \end{tabular}\par}
  \vspace*{1.05in}
\end{center}
\clearpage

% ------------------ Revision History ------------------
\frontmattertitle{Revision History}
\vspace{7pt}
\begin{tabularx}{\textwidth}{|L{2.10cm}|L{1.30cm}|Y|L{2.60cm}|}
  \hline
  \textbf{Date} & \textbf{Version} & \textbf{Description} & \textbf{Author}\\ \hline
  \upd{28 Nov 2026} & \upd{1.0} & \upd{New document in PA4: system test plan for the release produced in this iteration. Test cases and results are kept in the instructor's workbooks PA4-Group02-Test-Cases.xlsx and PA4-Group02-Test-Report.xlsx.} & \upd{Group 02}\\ \hline
\end{tabularx}

\medskip
\upd{\textit{Convention: in all PA4 documents, text written in blue was added or changed in PA4.}}
\clearpage

\tableofcontents
\clearpage

\frontmattertitle{\DocumentTitle}

% =====================================================================
\section{Introduction}
\subsection{Purpose}
This Test Plan describes how Group 02 performs system testing of the Medication Safety Assistant
release built in the PA4 iteration. It defines the features to be tested, the testing approach, the
environment, the pass/fail criteria, the schedule and the responsibilities, so that the results reported
in the Test Report can be traced back to the use cases of the requirements stage.

\subsection{Scope}
The plan covers \textbf{functional system testing} (black-box) of the twenty use cases UC-01 to UC-20
documented in the PA2/PA3 Use-Case Specification. Use cases whose main flow is implemented end to end
(server and user interface) are executed in this iteration; use cases that exist only as UI prototypes
or are reserved for future work are listed with the reason they are not fully executed.

\subsection{References}
\begin{itemize}
  \item PA2/PA3 Use-Case Specification (UC-01 to UC-20) and PA3 Software Architecture Document.
  \item PA4 Software Architecture Document v2.0 (deployment and implementation views).
  \item PA4-Group02-Test-Cases.xlsx and PA4-Group02-Test-Report.xlsx (instructor templates v1.1).
  \item Source repository: \url{https://github.com/hyt2434/Medical_Safety_Assistant}.
\end{itemize}

% =====================================================================
\section{Test Items}
\begin{longtable}{|L{4.2cm}|L{11.3cm}|}
\hline \textbf{Item} & \textbf{Description}\\ \hline \endhead
Web client & React/TypeScript single-page application (\texttt{apps/frontend}), production build.\\ \hline
Application server & FastAPI REST API (\texttt{apps/backend}) with routes for authentication, administration, intake/OCR, reference data, HCP review (Module 4) and reminders (Module 5).\\ \hline
Domain modules & \texttt{src/module1\_intake} (manual, voice, OCR, chatbot), \texttt{module2\_medications}, \texttt{module3\_interactions}, \texttt{module4\_hcp}, \texttt{module5\_reminders}, \texttt{shared} (auth, database repositories).\\ \hline
Database & PostgreSQL 18 schemas \texttt{identity}, \texttt{msa}, \texttt{msa\_ref} created by migrations 001--005 with row-level security and domain guards.\\ \hline
Release under test & Branch \texttt{dev} at commit \texttt{de32207c} (PR \#86), tagged as release \texttt{v0.4.0} on \texttt{main}.\\ \hline
\end{longtable}

% =====================================================================
\section{Features to Be Tested}
Each feature corresponds to one use case. At least five functional test cases are written for every use case
(main flow, invalid input, boundary values, permission/role checks, and error or recovery flows).
\begin{longtable}{|L{1.35cm}|L{5.4cm}|L{4.1cm}|L{4.2cm}|}
\hline \textbf{ID} & \textbf{Feature (use case)} & \textbf{Implementation level} & \textbf{Execution in PA4}\\ \hline \endhead
UC-01 & Sign In & Server + UI & Executed (automated)\\ \hline
UC-02 & Enter and Confirm Medication Information & Server + UI & Executed (automated)\\ \hline
UC-03 & Check Medication Safety & Server + UI & Partly; DDI cases need local reference data\\ \hline
UC-04 & Arrange HCP Review & Server + UI (new in PA4) & Executed (automated)\\ \hline
UC-05 & Review Medication Safety Case & Server + UI (new in PA4) & Executed (automated)\\ \hline
UC-06 & Set Up Medication Reminders & Server + UI (new in PA4) & Executed (automated)\\ \hline
UC-07 & Manage Medication Reminders & Server + UI (new in PA4) & Executed; known open defects\\ \hline
UC-08 & Record Dose Status & Server + UI (new in PA4) & Executed (automated)\\ \hline
UC-09 & View Adherence History & Server + UI (new in PA4) & Executed (automated)\\ \hline
UC-10 & View Medication and Safety History & UI prototype & Partly\\ \hline
UC-11 & Ask a Medication Question & Server + UI & Partly; model-dependent cases opt-in\\ \hline
UC-12 & Generate Patient Health Summary & UI prototype & Not executed (no server API)\\ \hline
UC-13 & Share Patient Health Summary with HCP & UI prototype & UI level only\\ \hline
UC-14 & Review Medication Reference Data & UI prototype & UI level only\\ \hline
UC-15 & Manage User Accounts and Roles & Server + UI & Executed (automated)\\ \hline
UC-16 & Report Technical Issue & UI prototype & UI level only\\ \hline
UC-17 & Handle Technical Support Request & UI prototype & UI level only\\ \hline
UC-18 & Manage Patient Health Profile & UI prototype & UI level only\\ \hline
UC-19 & Manage Personal Medication Cabinet & Future development & Not executed\\ \hline
UC-20 & Manage Personal Account & Server (view) + UI & Partly\\ \hline
\end{longtable}

\section{Features Not to Be Tested}
\begin{itemize}
  \item Clinical correctness of medical content: the product is a research prototype and does not replace
        professional judgement; only the system behaviour defined in the use cases is tested.
  \item Performance, load and penetration testing beyond the existing automated security checks.
  \item UC-19 (reserved for future development) and push/e-mail notification delivery (reminders are in-app only).
  \item Real-model accuracy of the local language model (Ollama/Qwen) and speech recognition (Whisper), which
        are opt-in and depend on hardware.
\end{itemize}

% =====================================================================
\section{Approach}
\subsection{Test level and type}
System testing of the integrated release through its public interfaces (browser UI and REST API),
using functional black-box techniques: equivalence partitioning, boundary-value analysis (times, dates,
time zones, dose windows), role/permission testing (RBAC), and alternative/exception flows from the
use-case specifications.

\subsection{Execution methods}
\begin{itemize}
  \item \textbf{Automated end-to-end UI tests} with Playwright (Chromium) at desktop (1280\,px) and mobile (375\,px) widths.
  \item \textbf{Automated API and integration tests} with pytest, including a full Module 4/5 flow on an isolated,
        marked PostgreSQL test database.
  \item \textbf{Manual exploratory tests} for cases without automation (recorded as \emph{Not run} until executed).
  \item An AI coding assistant was used to help draft test cases from the use-case specifications; every result
        recorded as Pass is backed by an executed automated test named in the test-case workbook.
\end{itemize}

\subsection{Regression}
The full automated suite is executed before every merge to \texttt{dev}; results are compared with the
previous baseline (27 Nov 2026) to detect regressions.

% =====================================================================
\section{Item Pass/Fail Criteria}
\begin{itemize}
  \item A test case \textbf{passes} when the actual result matches the expected result in every step.
  \item A test case \textbf{fails} when any step deviates; a defect is recorded in the Bug report.
  \item The release is \textbf{accepted} when the main flows of at least two main use-case groups pass,
        no executed test case of severity High is failing, and all open defects are documented with a plan.
\end{itemize}

\section{Suspension Criteria and Resumption Requirements}
Testing is suspended when the build fails, the database cannot be migrated, or sign-in (UC-01) fails,
because these block all other features. Testing resumes after the blocking defect is fixed and the
smoke tests (sign-in, medication intake, reminders) pass again.

\section{Test Deliverables}
Test Plan (this document); Test-case workbook; Test Report workbook (summary, bug report, screenshots);
automated test scripts in \texttt{tests/} and \texttt{apps/frontend/tests/}; execution logs.

% =====================================================================
\section{Environmental Needs}
\begin{longtable}{|L{3.6cm}|L{11.9cm}|}
\hline \textbf{Component} & \textbf{Specification}\\ \hline \endhead
Client devices & Windows 11 laptop; desktop browser (Chromium/Microsoft Edge, Firefox); mobile viewport 375\,px.\\ \hline
Server & Python 3.14, FastAPI/Uvicorn, Node.js for the frontend build and OCR helper.\\ \hline
Database & PostgreSQL 18; separate marked test database for destructive integration tests.\\ \hline
Optional services & Ollama with a local Qwen model (chatbot), faster-whisper (voice), Tesseract (OCR).\\ \hline
Test data & Synthetic patient accounts and synthetic medication records only; no real patient data.\\ \hline
Tools & pytest, Playwright, Git/GitHub (pull requests, releases).\\ \hline
\end{longtable}

\section{Responsibilities}
\begin{longtable}{|L{3.8cm}|L{11.7cm}|}
\hline \textbf{Member} & \textbf{Responsibility}\\ \hline \endhead
Hoàng Yến Trang & Test lead; Modules 4 and 5, User Management (UC-01, UC-15, UC-20) and OCR intake; test plan and report.\\ \hline
Cao Hoàng Long & Modules 2 and 3 and the database; safety checks (UC-03) and reference data (UC-14); provides reference datasets.\\ \hline
Lỗ Chí Thành & Voice and manual input (UC-02), reminders review (UC-06--UC-09).\\ \hline
\end{longtable}

\section{Schedule}
\begin{longtable}{|L{3.2cm}|L{9.0cm}|L{3.3cm}|}
\hline \textbf{Date} & \textbf{Activity} & \textbf{Owner}\\ \hline \endhead
27 Nov 2026 & Baseline run on \texttt{dev}; defects B001--B004 recorded & Trang\\ \hline
27--28 Nov 2026 & Modules 4/5 implemented; code reviews and fixes (B008--B013) & Trang\\ \hline
28 Nov 2026 & Full system run, test-case execution, report & Trang\\ \hline
29 Nov 2026 & Release v0.4.0 and PA4 submission & Group 02\\ \hline
Next iteration & Manual execution of Not-run cases; fixes for B005--B007, B014 & Group 02\\ \hline
\end{longtable}

\section{Risks and Contingencies}
\begin{longtable}{|L{5.5cm}|L{10.0cm}|}
\hline \textbf{Risk} & \textbf{Contingency}\\ \hline \endhead
Reference datasets missing on a test machine (B001) & Share the dataset package from the database owner; mark affected cases Not run instead of Pass.\\ \hline
Local AI models unavailable or slow & Keep model-dependent tests opt-in; test the non-model behaviour (limits, no automatic record changes).\\ \hline
Schema change needed for schedule proposals (B005) & Plan migration 006 with the database owner; keep the current verified schedule behaviour documented.\\ \hline
Limited time before the deadline & Prioritise the main flows of UC-01--UC-09; record remaining cases as Not run with reasons.\\ \hline
\end{longtable}

\section{Approvals}
\begin{tabularx}{\textwidth}{|Y|L{4cm}|L{3cm}|}
\hline \textbf{Name} & \textbf{Role} & \textbf{Date}\\ \hline
Hoàng Yến Trang & Test lead & 28 Nov 2026\\ \hline
Cao Hoàng Long & Team member & \\ \hline
Lỗ Chí Thành & Team member & \\ \hline
\end{tabularx}

\end{document}
